\BourbakiExercisesSection

\begin{enumerate}
\def\labelenumi{\arabic{enumi}.}
\tightlist
\item
  设 A 为整环，K 为其分式域。
\end{enumerate}

\begin{enumerate}
\def\labelenumi{(\alph{enumi})}
\item
  证明外幂 \(\bigwedge^{2}\mathbf{K}\) （即 A-模 \(\mathbf{K}\) 的外幂）缩减为 0。
\item
  对每个 A-模 \(\mathbf{E} \subset \mathbf{K}\) ，证明每个外幂 \(\bigwedge \mathbf{E}\) （\(p \geqslant 2\)）都是挠 A-模。
\item
  若 E 为 II，§7，习题 31 中定义的 A-模，证明 \(\bigwedge\) E 是非零的。给出一个整环 A 和包含在 A 的分式域 K 中的 A-模 E 的例子，使得任何外幂 \(\bigwedge^{p}\) E 都不化为 0。
\end{enumerate}

\begin{enumerate}
\def\labelenumi{\arabic{enumi}.}
\setcounter{enumi}{1}
\item
  对于环 A 的每个理想 a 和每个 A-模 E，外代数 \(\bigwedge(\mathrm{E}/\mathfrak{a}\mathrm{E})\) 同构于 \((\bigwedge\mathrm{E})/\mathfrak{a}(\bigwedge\mathrm{E})\) 。
\item
  给出一个具有如下性质的环 A 的例子：在 A-模 \(E = A^{2}\) 中，存在一个子模 F，使得，若 \(j: F \to E\) 是典范单射，则 \(\bigwedge^{2} j\) 恒为零，即使 \(\bigwedge^{2} E\) 与 \(\bigwedge^{2} F\) 都不化为 0（II，§4，习题 2）。
\item
  设 E 是维数为 \(n \geqslant 3\) 的自由 A-模。证明：若 \(l \leqslant p \leqslant n - 1\) ，则 A-模 \(\bigwedge^{p}E\) 与 \(\bigwedge^{n-p}E\) 同构；但若 \(2p \neq n\) ，则不存在同构 \(\phi\) （从 \(\bigwedge^{p}E\) 到 \(\bigwedge^{n-p}E\) 上），使得对 E 的每个自同构 u 都有 \(\phi \circ \left(\bigwedge^{p}u\right) = \left(\bigwedge^{n-p}u\right) \circ \phi\) （若 \((e_{i})\) 是 E 的一组基，考虑满足 \(u(e_{i}) = e_{i} + e_{j}, y(e_{k}) = e_{k}\) （对 \(k \neq i\) ）的自同构 u，并让 i 和 j 取遍所有可能的值）。
\item
  设 K 是一个交换域，E、F 是 K 上的两个向量空间，u 是从 E 到 F 的秩为 r 的线性映射。证明：若 \(p \leqslant r\) ，则 \(\bigwedge^{p}u\) 的秩等于 \(\binom{r}{p}\) ；且若 p \textgreater{} r，则 \(\bigwedge^{p}u\) 恒为零（在 E 中取一组基，使其中 r 个向量构成与 \(u^{-1}(0)\) 互补的一个子空间的基，而其余向量构成 \(u^{-1}(0)\) 的基）。
\end{enumerate}

¶ 6. 设 K 为一个交换域，E 为 K 上的向量空间。

\begin{enumerate}
\def\labelenumi{(\alph{enumi})}
\item
  为使元素 \(z = \sum_{p=0}^{\infty} z_p \left( \text{其中 } z_p \in \bigwedge^p \mathbf{E} \right)\) （属于外代数 \(\bigwedge \mathbf{E}\) ，即 \(\mathbf{E}\) 的外代数）可逆，必须且只需 \(z_0 \neq 0\) （先在 \(\mathbf{E}\) 为有限维时证明，然后过渡到一般情形，并注意到对所有 \(z \in \bigwedge \mathbf{E}\) ，存在一个有限维子空间 \(\mathbf{F}\) （包含于 \(\mathbf{E}\)），使得 \(z \in \bigwedge \mathbf{F}\) ）。
\item
  设 K 满足在 K 中 \(2 \neq 0\) 。证明：若 E 是无限维的或有限维且维数为偶数，则代数 \(\Lambda E\) 的中心由这样的元素组成：对所有奇数指标 p，\(z_{p} = 0\) 。若 E 的维数 n 为奇数，则 \(\Lambda E\) 的中心是上述子空间与 \(\Lambda^{n}E\) 之和。在所有情形下， \(\Lambda E\) 的中心都与 \(\Lambda E\) 中与 \(\Lambda E\)
\end{enumerate}

\begin{enumerate}
\def\labelenumi{\arabic{enumi}.}
\setcounter{enumi}{6}
\item
  的所有可逆元素可交换的元素所成的集合相同。对每个 A-模 E，证明：若我们记 \([u, v] = u \wedge v - v \wedge u\) （对 \(\bigwedge E\) 的任意两个元素），则对任意三个元素 \(u, v, w\) （属于 \(\bigwedge E\)），有 \([(u, v], w] = 0\) 。
\item
  设 E 是一个自由 A-模。
\end{enumerate}

\begin{enumerate}
\def\labelenumi{(\alph{enumi})}
\tightlist
\item
  证明在 \(\mathbf{T}^n (\mathbf{E})\) 中， \(\mathfrak{J}_n''\) （第 1 号）是一个自由子 A-模，等于自同态 \(z\mapsto a.z\) 的核，并且有一个补，该补是自由 A-
\end{enumerate}

模；因此 \(\bigwedge E\) 典范同构于 \(\mathsf{A}_n^{\prime \prime}(\mathbf{E})\) 。若此外在 A 中方程 \(2\xi = 0\) 蕴含 \(\xi = 0\) ，则 \(\mathsf{A}_n^{\prime \prime}(\mathbf{E}) = \mathsf{A}_n^{\prime}(\mathbf{E})\) 。

*(b) 设 A 是特征为 p \textgreater{} 0 的域，且 \(p \neq 2\) 。则 \(\mathbf{A}_{p}^{\prime}(\mathbf{E}) = \mathbf{A}_{p}^{\prime\prime}(\mathbf{E})\) 在 \(\bigwedge^{p} E\) 中的典范像为零。

\begin{enumerate}
\def\labelenumi{(\alph{enumi})}
\setcounter{enumi}{2}
\tightlist
\item
  设 \(\mathbf{A}\) 是特征为 2 的域。则 \(\mathbf{A}_n^\prime (\mathbf{E}) = \mathbf{S}_n^\prime (\mathbf{E}),\) \(\mathbf{A}_n''(\mathbf{E}) = \mathbf{S}_n''(\mathbf{E})\) ，但一般地 \(\mathbf{A}_n^\prime (\mathbf{E})\neq \mathbf{A}_n''(\mathbf{E})\cdot *\)
\end{enumerate}

\begin{enumerate}
\def\labelenumi{\arabic{enumi}.}
\setcounter{enumi}{8}
\tightlist
\item
  设 E 是一个 A-模。从 \(\mathsf{T}^{n}(\mathbf{E})\) 到一个 A-模 F 中的线性映射 g 的斜对称化，按定义是从 \(\mathsf{T}^{n}(\mathbf{E})\) 到 F 中的线性映射 ag，它使得 \(ag(z) = g(\boldsymbol{a}z)\) 对所有 \(z \in \mathsf{T}^{n}(\mathbf{E})\) 成立；若 \(\phi: E^{n} \to T^{n}(\mathbf{E})\) 是典范映射，则 n-线性映射 \(ag \circ \phi\) 也称为 n-线性映射 \(g \circ \phi\) 的斜对称化。由习题 8(a) 推出，若 E 是自由的，则每个从 \(E^{n}\) 到 F 中的交错 n-线性映射都是某个从 \(E^{n}\) 到 F 中的 n-线性映射的斜对称化。
\end{enumerate}

¶ 10. 设 E 是一个具有 n 个元素的生成系的 A-模。证明每个从 \(E^{n}\) 到一个 A-模 F 中的交错 n-线性映射都是某个从 \(E^{n}\) 到 F 中的 n-线性映射的斜对称化。（设 \(a_{j}\) ( \(1 \leqslant j \leqslant n\) ) 为 E 的生成元，且 \(\phi: A^{n} \rightarrow E\) 为满足 \(\phi(e_{j}) = a_{j}\) （对所有 j）的同态，其中 \((e_{j})\) 是 \(A^{n}\) 的典范基。若 \(f: E^{n} \rightarrow F\) 是交错 n-线性映射，则 \(f_{0}: (x_{1}, \ldots, x_{n}) \mapsto f(\phi(x_{1}), \ldots, \phi(x_{n}))\) 也是；证明：若 \(g_{0}\) 是从 \((A^{n})^{n}\) 到 F 中的 n-线性映射，使得 \(f_{0} = a g_{0}\) ，则 \(g_{0}\) 也可以写成 \((x_{1}, \ldots, x_{n}) \mapsto g(\phi(x_{1}), \ldots, \phi(x_{n}))\) ，其中 g 是从 \(E^{n}\) 到 F 中的 n-线性映射。）

¶ 11. 设 K 是一个交换域，其中 \(2 = 0\) ，A 为多项式环 \(\mathrm{K}[\mathrm{X}, \mathrm{Y}, \mathrm{Z}]\) ，E 为 A 模 \(\mathrm{A}^3/\mathrm{M}\) ，其中 M 是 \(\mathrm{A}^3\) 的由元素 \((\mathrm{X}, \mathrm{Y}, \mathrm{Z})\) （属于 \(\mathrm{A}^3\)）生成的子模；最后，设 F 为 A 关于理想 \(\mathrm{AX}^2 + \mathrm{AY}^2 + \mathrm{AZ}^2\) 的商 A-模。证明存在一个从 \(\mathrm{E}^2\) 到 F 中的交错双线性映射，它不是任何从 \(\mathrm{E}^2\) 到 F 中的双线性映射的斜对称化，但在 \(\mathsf{T}^2(\mathbf{E})\) 中模 \(\mathfrak{S}_2''\) 等于自同态 \(z \mapsto az\) 的核（把 \(\mathsf{T}^2(\mathbf{E})\) 看作 \(\mathsf{T}^2 (\mathbf{A}^3)\) 的商模，并证明 \(z\mapsto az\) 的核是对称张量模 \(\mathsf{S}_2^{\prime}(\mathbf{A}^3)\) 的典范像）。

¶ 12. 在习题11的记号下，设B为商环

\[
\mathrm{A} / (\mathrm{AX} ^ {2} + \mathrm{AY} ^ {2} + \mathrm{AZ} ^ {2})
\]

并设 G 为 B-模 E ⊗A B。证明在 T²(G) 中，模 \(J_{2}^{\prime\prime}\) 不同于自同态 \(z \mapsto az\) 的核（方法类似于习题 11 的方法）。

\begin{enumerate}
\def\labelenumi{\arabic{enumi}.}
\setcounter{enumi}{12}
\tightlist
\item
  \begin{enumerate}
  \def\labelenumii{(\alph{enumii})}
  \tightlist
  \item
    设 M 是 N 型分次 A-模。证明在代数 \(S(M)\) （相应地 \(\bigwedge(M)\)）上存在且仅存在一种 N 型分次，使得 \(S(M)\) （相应地 \(\bigwedge(M)\)）成为分次 A-代数，并在 M 上诱导给定的分次。证明：若 M 中偶数次齐次元素全为零（此时 M 上的分次称为奇分次），则如此得到的分次代数 \(\bigwedge(M)\) 是交错的（ \(\S 4\) ，第 9 号，定义 7）。
  \end{enumerate}
\end{enumerate}

\begin{enumerate}
\def\labelenumi{(\alph{enumi})}
\setcounter{enumi}{1}
\tightlist
\item
  给定一个 \(\mathbf{N}\) 型分次 A-模，考虑以下泛映射问题： \(\Sigma\) 是反交换分次代数结构（§ 4，第 9 号，定义 7）的种，态射是分次代数的 A-同态（§ 3，第 1 号），而 \(\alpha\) -映射是 M 到反交换分次 A-代数中的次数为 0 的分次同态。设 \(\mathbf{M}^{-}\) （相应地 \(\mathbf{M}^{+}\)）表示 M 的由奇（相应地，偶）次数齐次元素生成的分次子模。证明分次代数 \(\mathbf{G}(\mathbf{M}) = \bigwedge (\mathbf{M}^{-}) \otimes_{\mathbf{A}} \mathbf{S}(\mathbf{M}^{+})\) 是反交换的；M 到这个代数中的一个典范单射 \(j\) 由写出 \(j(x) = x \otimes 1\) 若 \(x \in \mathbf{M}^{-}, j(x) = 1 \otimes x\) 若 \(x \in \mathbf{M}^{+}\) 来定义。证明分次代数 \(\mathbf{G}(\mathbf{M})\) 与单射 \(j\) 构成上述泛映射问题的一个解。
\end{enumerate}

\[
\mathbf {G} (\mathbf {M}) = \bigwedge (\mathbf {M} ^ {-}) \otimes_ {\mathbf {A}} \mathbf {S} (\mathbf {M} ^ {+})
\]

称为分次A-模M上的泛反交换代数。

\begin{enumerate}
\def\labelenumi{(\alph{enumi})}
\setcounter{enumi}{2}
\tightlist
\item
  证明关于泛反交换代数的命题 2、3 和 4 的类似结论。考虑 M 允许有由齐次元素组成的一组基的情形：在此情形下，明确给出泛反交换代数的一组基。
\end{enumerate}

\begin{enumerate}
\def\labelenumi{\arabic{enumi}.}
\setcounter{enumi}{13}
\tightlist
\item
  设 M 是一个投射 A-模，且令 \(x_{1}, \ldots, x_{n}\) 是 M 的线性无关元素；对每个子集 \(H = \{i_{1}, \ldots, i_{m}\}\) ——由 \(m \leqslant n\) 个区间 \([1, n]\) 的元素组成——其中 \(i_{1} < i_{2} < \cdots < i_{m}\) ，我们记
\end{enumerate}

\[
x _ {\mathrm{H}} = x _ {i _ {1}} \wedge x _ {i _ {2}} \wedge \dots \wedge x _ {i _ {m}}.
\]

证明当 H 遍历 \([1, n]\) 的具有 m 个元素的子集时，诸 \(x_{H}\) 在 \(\bigwedge(M)\) 中线性无关。

\begin{enumerate}
\def\labelenumi{\arabic{enumi}.}
\setcounter{enumi}{14}
\tightlist
\item
  设 M 是一个 n 维自由 A-模。证明 M 的每个具有 n 个元素的生成系都是 M 的一组基。
\end{enumerate}

